Human--wildlife conflict (HWC) -- interactions between humans and wildlife that
produce negative outcomes such as crop damage, livestock loss, and threats to
human safety -- is one of the most pressing challenges in contemporary wildlife
conservation \citep{Konig2020, Madden2004, Zimmermann2020}. Beyond its ecological
consequences, HWC is fundamentally a social issue: whether it escalates or
subsides depends heavily on how the people who bear its costs come to feel about
the wildlife causing them. Within HWC research, \emph{tolerance} -- the ability
and willingness of residents to absorb the real or perceived costs of living
alongside wildlife -- has emerged as the key indicator of whether coexistence is
achievable \citep{ExpositoGranados2019, KanskyKidd2024}. Tolerance is shaped not
only by tangible costs and benefits such as economic losses or compensation
payments, but also by intangible dimensions -- emotional responses, perceived
risk, and the social meanings people attach to wildlife
\citep{KanskyKnight2014}.

Compensation schemes, which reimburse residents for wildlife-caused damage, are
a widely used policy instrument intended to offset these costs and, in doing so,
raise tolerance \citep{Hemminger2025, Kichloo2023}. In China, wild boar (\emph{Sus
scrofa}) populations have expanded rapidly in recent decades, and with them the
frequency of crop damage and compensation claims in and around national parks.
Whether these compensation payments actually achieve their intended social
effect -- higher tolerance among the people who receive them -- is the open
empirical question this project addresses. Preliminary evidence from our project
partner's fieldwork suggested the opposite of the policy assumption:
compensation recipients did not show higher tolerance than non-recipients, and
if anything the association ran the other way. Understanding whether that
pattern holds up under careful statistical scrutiny, and if so, why compensation
might fail to translate into tolerance, is central to designing socially
sustainable wildlife-management policy.

This report documents a statistical consulting collaboration between the
Department of Statistics at Ludwig-Maximilians-Universit\"at M\"unchen and
project partner Mengxi Kou (Rachel Carson Center for Environment and Society),
built around household survey data from two counties bordering Giant Panda
National Park, Sichuan Province, China. The analysis is organized around three
research questions, unchanged from the original project proposal:

\begin{enumerate}
  \item[\textbf{RQ1}] Does compensation for wild boar damage have a direct
    effect on tolerance toward wild boar, or is its influence mediated by other
    social and psychological factors?
  \item[\textbf{RQ2}] What roles do tangible costs, intangible costs, and
    intangible benefits play in shaping tolerance, and how do these factors
    interact with compensation experience?
  \item[\textbf{RQ3}] How robust are these relationships under different model
    specifications and across groups with and without compensation experience?
\end{enumerate}

Section~\ref{data} describes the survey and the data-quality checks that shaped
every later modelling decision, most importantly a strong confound between
compensation status and county of residence. Section~\ref{measurement} covers
the measurement work: which survey items reliably capture the constructs the
original Wildlife Tolerance Model \citep{Kansky2016} assumes, and which do not.
Section~\ref{model} presents the main structural model and the findings for
RQ1 and RQ2. Section~\ref{robustness} subjects those findings to two independent
robustness exercises -- a specification-vibration analysis and a permutation-based
multi-group comparison -- addressing RQ3. Section~\ref{discussion} discusses what
the findings do and do not support, and what would be needed to strengthen the
next round of data collection. Section~\ref{conclusion} summarizes the project's
conclusions.

A note on terminology: throughout this report we use short construct codes
(e.g.\ \emph{IC} for Intangible Cost, \emph{CE0} for compensation participation)
consistently with the analysis scripts and the presentation delivered to the
project partner, because these codes make cross-referencing tables and models
unambiguous. Every code is spelled out in full at first use in the running text,
and Appendix~\ref{app:glossary} collects all of them in a single glossary table
for quick reference.
